\section{Out-of-Distribution Transfer: EDF-S}
\label{sec:ood}

% --- motivation ---
The value proposition of a foundation model is transfer without retraining. All
results above are in the ECDFS field on which the downstream heads were tuned.
Here we apply the frozen \jaisp{} stack --- foundation, detector, and
joint-centroid astrometry --- to 72 paired tiles in the independent EDF-S field
(tract 2394) with no retraining or refitting.

% --- caveat on this pass ---
\todo{State the final data condition. The first pass used \textit{Euclid} tiles
without variance maps, requiring a median-absolute-deviation RMS fallback for
VIS/NISP (Rubin variance was present and used directly). Results below are
either (a) to be reported on real-variance tiles, or (b) reported with the
fallback and the real-variance rerun as a robustness check.}

% --- what transfers ---
Detection transfers cleanly (456 sources per tile versus 413 on ECDFS; no
processing errors), and \textit{Euclid} NISP astrometry matches the
in-distribution residual ($35$--$41\mas$ versus $40\mas$). The signal-to-noise
scaling of the residual is reproduced. \todo{Final Rubin optical
($g\,r\,i\,z$) transfer result: the first (fallback) pass showed a
$\sim$1.5$\times$ larger matched-SNR residual whose origin --- the RMS fallback
perturbing the joint-fit weighting versus genuine field differences --- is
resolved by the real-variance rerun. Report the resolved outcome here.}

\begin{figure}
  \centering
  % \includegraphics[width=\linewidth]{ood_astrometry_compare.png}
  \caption{Per-band and SNR-binned astrometric residual, EDF-S (OOD) versus
    ECDFS. \todo{final figure from notebook 15.}}
  \label{fig:ood}
\end{figure}

% --- detection completeness OOD ---
The bright-star detection suppression (\S\ref{sec:results:detection})
reproduces and is somewhat stronger in EDF-S ($\mathrm{corr}=-0.71$),
consistent with its origin in the input normalization rather than the field.
\todo{MER-catalog completeness in EDF-S.}
